problem:  
Let $(M^{2n+1}, g, \xi_E, \eta, \Phi)$ be a compact Sasaki–Einstein manifold with Reeb vector field $\xi_E$, and fix its underlying CR structure $(D, J)$.  
The Reeb cone $\mathcal{C}_R$ parametrizes all Reeb vector fields compatible with $(D, J)$.  
For $\xi' \in \mathcal{C}_R$, let $V(\xi')$ denote the total volume of the corresponding normalized Sasakian structure.  
By Martelli–Sparks–Yau, $\xi_E$ is the unique global minimizer of $V$ under a normalization constraint.

Let $\Delta_B$ denote the transverse basic Laplacian acting on basic functions and $\mathcal{L}_B$ the transverse Lichnerowicz operator acting on basic $(1,1)$-forms, both defined with respect to $\xi_E$. Their spectra are discrete and nonnegative.  

**Open Problem (spectral–stability formulation):**  
Assume that $\xi_E$ is a strict local minimizer of the normalized volume functional in $\mathcal{C}_R$, with nondegenerate Hessian $H_V(\xi_E)$.  
Is there a *quantitative relation* between $H_V(\xi_E)$ and the lowest nonzero eigenvalues of $\Delta_B$ or $\mathcal{L}_B$?  
In particular, does there exist a function $\Psi$, depending only on the CR structure $(D,J)$ and the dimension $n$, such that
\[
\lambda_1(\mathcal{L}_B) \geq \Psi(H_V(\xi_E))
\]
where $\lambda_1(\mathcal{L}_B)$ is the first nonzero eigenvalue of $\mathcal{L}_B$?  
Equivalently, do variational properties of the volume functional in the Reeb cone govern the spectral gap controlling infinitesimal stability of the Sasaki–Einstein metric under transverse Kähler deformations of fixed CR type?

Why is it a "good" problem:  
This problem transforms the physical intuition of stability in AdS/CFT into a rigorous geometric analysis question: can convexity of the Reeb-volume functional dictate the spectral stability of the underlying transverse geometry? It avoids unverifiable universal constants and instead seeks a structural link between two deep notions—variational convexity in contact–complex geometry and analytic positivity in the spectra of natural geometric operators.  

If such a relationship exists, it would mirror and geometrically justify the stability principles underlying “a-maximization” in the physical dual field theory, replacing them with a differential-geometric inequality involving $H_V(\xi_E)$ and the Laplacian/Lichnerowicz spectrum. The statement is mathematically well-defined, conceptually simple, and profoundly unifying—bringing together convex analysis on the Reeb cone, spectral theory, and geometric stability of Sasaki–Einstein metrics. Proving or refuting it would illuminate the analytic foundations of Sasaki–Einstein moduli and potentially open a new bridge between geometric analysis and high-energy theory.